\chapter{Ordered collisions and cut reconstruction}
\label{cycuts:ch}
\begingroup
\def\DMod{\operatorname{DMod}}
\def\RHom{R\operatorname{Hom}}
\def\W{\mathcal W_X}
\def\C{\mathcal C_X}
\def\cut{\star_{\mathrm{cut}}}
\def\Bcut{B_{\mathrm{cut}}}
\def\Ocut{\Omega_{\mathrm{cut}}}

Consider three positions in the order $(x,y,x)$, where $x\ne y$.
Neither cut
\[
 (x)\mid(y,x),\qquad (x,y)\mid(x)
\]
separates disjoint sets of positions. The outer collision therefore
survives a construction that permits only consecutive cuts. Collapsing
the two occurrences of $x$ gives two distinct positions, which do admit
a cut. This discrepancy determines both a useful reconstruction theorem
and a limitation of its comparison with full collision geometry.

We begin with the cut product and its bar and cobar constructions.
Strong cut coproducts identify a finite localization complex; its support
contains every configuration whose equalities obstruct all cuts. We then
construct the ordered collision atlas and its product over arbitrary
ordered subsets. It has its own bar--cobar equivalence. For the dualizing
coalgebra, the map from its cobar to the cut cobar has a kernel with
explicit principal-part classes, including negative-degree classes in
arbitrarily large arity.

The last comparison concerns factorization. Invertibility of each
coproduct component differs from invertibility of their joint map.
Full exchange data identify the atlas with the corresponding word
systems, while a local field realization must also preserve the actual
collision operations. Throughout, these statements concern associative
algebras and coassociative coalgebras on specified geometric carriers.

\section{Coefficients on ordered words}
\label{cycuts:sec:words}

Let $k$ be a field of characteristic zero and let $X$ be a smooth
separated curve over $k$. Write $\DMod(Y)$ for the enhanced stable
category of unbounded right differential-module complexes on a smooth
scheme $Y$, with quasi-coherent underlying modules. External products,
direct images and extraordinary inverse images are derived. Our shift
convention is $M[r]^i=M^{i+r}$, with differential $(-1)^r d_M$.

An ordered word coefficient system consists of a complex $F_0$ of
$k$-vector spaces and objects $F_n\in\DMod(X^n)$ for $n>0$.
For every monotone surjection $\alpha:[n]\twoheadrightarrow[m]$, it
includes an equivalence
\[
 d_\alpha^!F_n\simeq F_m,
 \qquad d_\alpha(y)_i=y_{\alpha(i)},
\]
compatible with identities, composites and their higher coherences.
Here $[n]=\{1,\ldots,n\}$ in its usual order. Let $D(k)$ be the
derived category of $k$-vector spaces, and let $\mathsf{OrdSurj}_+$
have the nonempty words $[n]$ as objects and monotone surjections
as arrows. The category of coefficient systems is
\[
 \W=D(k)\times
 \varprojlim_{\mathsf{OrdSurj}_+}\DMod(X^n).
\]
The first factor is isolated: there are no collision maps joining it
to a nonempty word.

For $1\le p<n$, define
\begin{equation}\label{cycuts:eq:opens}
 U_p^{(n)}=\{x_i\ne x_j\text{ whenever }i\le p<j\},
 \qquad j_p:U_p^{(n)}\hookrightarrow X^n.
\end{equation}
Every pair crossing the cut occurs in this condition, including pairs
of nonadjacent labels. Set $U_0^{(n)}=U_n^{(n)}=X^n$. Define
\begin{equation}\label{cycuts:eq:product}
 (F\cut G)_n=
 \bigoplus_{p=0}^n j_{p,*}j_p^!(F_p\boxtimes G_{n-p}).
\end{equation}
At an empty end, the external product means tensoring by the scalar
complex $F_0$ or $G_0$. The candidate unit is
\[
 K_0=k,\qquad K_n=0\quad(n>0).
\]
This empty-support unit is distinct from a vacuum coefficient on $X$.

We will repeatedly use open base change. If $j:U\hookrightarrow Y$
is open, $f:Y'\to Y$, and $j':U'\hookrightarrow Y'$ is its pullback,
then
\begin{equation}\label{cycuts:eq:basechange}
 f^!j_*\simeq j'_*f_U^!.
\end{equation}
To see this, let $i,i'$ be the closed complements. Recollement gives
$i^!j_*=0$. The composite $fi'$ factors through $i$, so
$i'^!f^!j_*=0$. The localization triangle identifies $f^!j_*M$
with the direct image of its restriction to $U'$, which is $f_U^!M$.
All comparisons restrict to transitivity of extraordinary inverse image
on $U'$. Full faithfulness of $j'_*$ therefore supplies their composition
coherences as well.

A consecutive partition divides a word into nonempty intervals,
listed in their original order.

\begin{proposition}\label{cycuts:prop:product}
Formula~\eqref{cycuts:eq:product} makes $\W$ a monoidal category with
unit $K$. Its iterated products are finite sums over consecutive
partitions, localized where distinct blocks have disjoint positions.
\end{proposition}
\begin{proof}
A monotone collision kills a cut inside one of its fibres: that fibre
would have to meet both sides at distinct positions. Every other cut
is the inverse image of exactly one quotient cut. Equation
\eqref{cycuts:eq:basechange} identifies the surviving summands with
the quotient product, compatibly with further collisions.

Each parenthesized product imposes the same final consecutive partition
and the same cross-block inequalities. Successive open direct images
identify its coefficient with the external product on that common open.
These identifications give associativity and all its coherences. The
two empty-end summands give the unit laws.
\end{proof}

Call a coefficient system \emph{positive} if its empty coefficient
vanishes. More generally, let $\mathcal W_{X,\ge a}$ consist of systems vanishing
in arities below $a$. A nonzero summand of a product requires the labels
used by each factor, so
\begin{equation}\label{cycuts:eq:arity}
 \mathcal W_{X,\ge a}\cut\mathcal W_{X,\ge b}\subseteq\mathcal W_{X,\ge a+b}.
\end{equation}
In particular, $r$ positive factors have zero coefficient at $n<r$.
Every reduced word construction below is finite at each configuration.
This finiteness uses arity, with no cohomological boundedness assumption.

\section{The two associative constructions}
\label{cycuts:sec:bar}

Here $E_1$ means homotopy-coherent associative. Let $A$ be an augmented
$E_1$ algebra for $\cut$, with $A_0=k$ and
augmentation ideal $a=\operatorname{fib}(A\to K)$. Its unit splits
$A\simeq K\oplus a$, and $a$ is positive. Write $s=[1]$, so
$|sx|=|x|-1$. The normalized bar has underlying word object
\begin{equation}\label{cycuts:eq:bar}
 \Bcut A=K\oplus\bigoplus_{r\ge1}(sa)^{\cut r}.
\end{equation}
Its coproduct deconcatenates words, retaining both empty ends. On a
strict differential graded representative, its degree-one coderivation
is determined by
\begin{equation}\label{cycuts:eq:bar-signs}
 b_1(sx)=-s(dx),\qquad b_2(sx,sy)=(-1)^{|x|}s(xy).
\end{equation}
Insertion after a suspended prefix contributes the Koszul sign of that
prefix. Thus for closed degree-zero inputs,
\[
 b[x|y]=[xy],\qquad
 b[x|y|z]=[xy|z]-[x|yz].
\]
The second differential on the last expression is zero by associativity.
These words denote geometric coefficient summands; they need not be
global pure tensors. A coherent $E_1$ model retains the higher faces
of its normalized bar diagram.

Let $C$ instead be a coaugmented counital $E_1$ coalgebra with $C_0=k$,
and put $c=\operatorname{cofib}(K\to C)$. Its reduced coproduct is
$\bar\Delta:c\to c\cut c$. With $t=[-1]$, define
\begin{equation}\label{cycuts:eq:cobar}
 \Ocut C=K\oplus\bigoplus_{r\ge1}(tc)^{\cut r}.
\end{equation}
Multiplication concatenates consecutive block lists. On a strict model,
the differential is the derivation
\begin{equation}\label{cycuts:eq:cobar-signs}
 d(tc)=-t(d_Cc)+
 \sum_{\bar\Delta c=\sum c'\otimes c''}
             (-1)^{|c'|+1}(tc')(tc'').
\end{equation}
The notation for the reduced coproduct abbreviates its localized
coefficient map. Internal differentials cancel the mixed terms in
$d^2$; disjoint splits cancel by the odd derivation rule; the two
refinements of one block into three cancel by coassociativity.

\begin{theorem}\label{cycuts:thm:duality}
Bar and cobar are inverse equivalences between augmented $E_1$ cut
algebras with empty coefficient $k$ and coaugmented counital $E_1$
cut coalgebras with empty coefficient $k$. The equivalence includes
their derived mapping spaces and higher structure maps.
\end{theorem}
\begin{proof}
Use the enhanced bar and cobar, with normalized expressions
\eqref{cycuts:eq:bar}--\eqref{cycuts:eq:cobar}. Their natural
comparisons have strict-model forms
\[
 \Ocut\Bcut A\longrightarrow A,
 \qquad C\longrightarrow\Bcut\Ocut C.
\]
The first uses the universal twisting map: a one-factor bar word
desuspends to its generator, and longer generators map to zero. The second
uses all iterated reduced coproducts. Positivity makes every component
of these formulas finite.

Flatten a word in $\Ocut\Bcut A$ into $N$ original factors of $a$.
Filter by $N$. Multiplication of original factors lowers this number.
In the associated graded, each of its $N-1$ gaps has two states: it
is inside a bar block, or between cobar factors. Changing the state
is the identity on the same final localized coefficient, in adjacent
degrees. The gap complex is the signed tensor product of these
two-term identity complexes. For $N>1$, contracting the first gap
contracts the whole complex. At $N=0,1$, the comparison is the
identity on $K$ and the original factor.

For $\Bcut\Ocut C$, use the decreasing filtration by the number
of original factors of $c$. Splitting a factor raises that number.
Its associated graded has the same identity gap complexes, with
the two states interchanged. The comparison is again an equivalence
in lengths zero and one and has contractible higher pieces.

At arity $n$, both filtrations have length at most $n$ by
\eqref{cycuts:eq:arity}. Thus their associated-graded contractions
prove both comparisons to be equivalences. In coherent models, higher
products lower the first filtration and higher coproducts raise the
second; the associated-graded gap diagrams are unchanged. Normalized
realization and totalization retain their higher faces. Their natural
comparisons and triangle coherences therefore give inverse functors
on the enhanced categories, including mapping spaces.
\end{proof}

The bar is a coalgebra and the cobar is an algebra. An exchanged
associative algebra requires a coefficient duality as well. For a
contravariant equivalence $\sharp$ on an admissible coefficient
subcategory, the coefficient transposition of Chapter~\ref{cyjets:ch}
takes the following form for the cut product:
\[
 A^!_\sharp=(\Bcut A)^\sharp,
 \qquad
 F^\sharp\star_\sharp G^\sharp=(F\cut G)^\sharp.
\]
The transposed deconcatenation multiplies $A^!_\sharp$ for this
transported tensor. Its carrier and admissibility hypotheses remain
part of the construction. No Lie--commutative duality is used to
identify it with an algebra of fields.

\section{Which collisions the cut cobar retains}
\label{cycuts:sec:support}

Write $\Delta_{ij}\subset X^n$ for the equality locus $x_i=x_j$,
$\Delta_n:X\hookrightarrow X^n$ for the small diagonal, and
$\Delta=\Delta_2$. Define the reduced closed support
\begin{equation}\label{cycuts:eq:Z}
 Z_1=X,\qquad
 Z_n=X^n\setminus\bigcup_{p=1}^{n-1}U_p^{(n)}
     =\bigcap_{p=1}^{n-1}\ \bigcup_{i\le p<j}\Delta_{ij}.
\end{equation}
A tuple lies in $Z_n$ precisely when every cut has an equality
crossing it. The first examples are
\begin{equation}\label{cycuts:eq:small-Z}
 Z_2=\Delta_{12},\qquad Z_3=\Delta_{13},\qquad
 Z_4=\Delta_{14}\cup(\Delta_{13}\cap\Delta_{24}).
\end{equation}
For the four-point formula, equality of the outer positions blocks
every cut. If those positions differ, failure of the first and last
cuts leaves the possibilities $12,34$ or $13,24$. The first still
admits the middle cut unless the outer positions become equal.
The second blocks all three cuts. These are identities of reduced
supports, without a claim about scheme multiplicities.

A coaugmented cut coalgebra is \emph{strong} if each of its supplied
coproduct maps
\[
 C_n|_{U_p^{(n)}}\longrightarrow
       (C_p\boxtimes C_{n-p})|_{U_p^{(n)}}
\]
is an equivalence. For a reduced closed support $Z$ with inclusion $i$,
write $R\Gamma_ZM=i_*i^!M$. Recollement expresses it as the fibre
of restriction followed by direct image from the complementary open.

\begin{theorem}\label{cycuts:thm:support}
For a strong cut coalgebra, the reduced cobar has a natural equivalence
\begin{equation}\label{cycuts:eq:local-cohomology}
 \overline{\Ocut C}_n\simeq R\Gamma_{Z_n}C_n[-1].
\end{equation}
It commutes with monotone collisions. Its map to $C_n[-1]$ is the
length-one projection of the cobar.
\end{theorem}
\begin{proof}
A subset $S\subseteq[n-1]$ of gaps specifies a consecutive partition
$P_S$. Let $U_S=\bigcap_{p\in S}U_p^{(n)}$, with inclusion $j_S$,
and put $L_S=j_{S,*}j_S^!$; the empty subset gives the identity.
The iterated strong coproduct identifies the entire refinement diagram
of cobar coefficients with $S\mapsto L_SC_n$. Coassociativity
identifies its faces and higher homotopies, not merely its vertices.

The nonempty subsets form the finite \v Cech diagram of the union
of cut opens. Adding the empty vertex takes its total fibre and
therefore gives $R\Gamma_{Z_n}C_n$. A partition with $r$ blocks
has $r-1$ cuts. Its cobar term has $r$ desuspensions, producing
the additional shift $[-1]$ in \eqref{cycuts:eq:local-cohomology}.

For the signs, consolidate the block shifts by
\begin{equation}\label{cycuts:eq:consolidation}
 (tc_1)\otimes\cdots\otimes(tc_r)\longmapsto
 (-1)^{\sum_{j=1}^r(r-j)|c_j|}
       (c_1\otimes\cdots\otimes c_r)[-r].
\end{equation}
Splitting block $i$ into degrees $a,b$ changes this exponent by
$\sum_{j<i}|c_j|+a$. Adding the cobar splitting exponent gives
$i$ modulo two. Thus the face has sign $(-1)^i$, the negative
of the corresponding \v Cech insertion sign; the internal differential
has sign $(-1)^r$. These are precisely the shifted-cube signs.

A monotone collision deletes cuts inside its fibres and identifies
the other cuts with all quotient cuts. Open base change identifies
both augmented cubes. In particular $d_\alpha^{-1}Z_n=Z_m$.
This proves compatibility, naturality and the assertion about the
empty vertex of the cube.
\end{proof}

Write $L_p=L_{\{p\}}$ and $L_{12}=L_{\{1,2\}}$. At two positions,
the resolution is
\[
 \operatorname{Tot}[C_2[-1]\xrightarrow{-\mathrm{loc}}L_1C_2[-2]].
\]
At three positions it is
\begin{equation}\label{cycuts:eq:three-cube}
 \operatorname{Tot}\left[
 C_3[-1]\xrightarrow{(-\mathrm{loc},-\mathrm{loc})}
 (L_1C_3\oplus L_2C_3)[-2]
 \xrightarrow{(u,v)\mapsto u|_{12}-v|_{12}}L_{12}C_3[-3]
 \right].
\end{equation}
The two composites cancel. For a smooth $Y$, let $p_Y:Y\to\operatorname{Spec}k$
be its structure map, $\omega_Y$ its canonical right differential
module, and $J_Y=p_Y^!k$. The dualizing word system is
\begin{equation}\label{cycuts:eq:J-word}
 J_n=p_{X^n}^!k=\omega_{X^n}[n],\qquad J_0=k,
\end{equation}
the coproduct is localization and is strong. Transitivity of $!$-pullback
gives its collision identifications. If $i_{13}(x,y)=(x,y,x)$, then
\begin{equation}\label{cycuts:eq:J-support}
 \overline{\Ocut J}_2\simeq\Delta_{12,*}J_X[-1],
 \qquad
 \overline{\Ocut J}_3\simeq i_{13,*}J_{X^2}[-1].
\end{equation}
The latter coefficient remains nonzero at $x_1=x_3\ne x_2$.

\subsection{The product on the localization complex}

The support equivalence retains an algebra structure through the
specified strong coproduct. Its inverse on $U_p^{(p+q)}$ is a map
\[
 m_p:(C_p\boxtimes C_q)|_{U_p}\longrightarrow C_{p+q}|_{U_p}.
\]
Write the finite resolution as
\[
 \mathcal K(C)_n=
 \operatorname{Tot}\!\bigoplus_{S\subseteq[n-1]}L_SC_n[-|S|-1].
\]
The internal differential on its $S$ term is $(-1)^{|S|+1}d$;
adding a gap $p$ has sign $(-1)^{1+\#\{a\in S:a<p\}}$.
For $S\subseteq[p-1]$ and $T\subseteq[q-1]$, put
\[
 r=|S|+1,\quad s=|T|+1,\quad
 G=S\cup\{p\}\cup(p+T).
\]
Their product is defined on $U_G$ by the restriction of $m_p$ and
then extended by $j_{G,*}$. On unshifted homogeneous coefficients,
\begin{equation}\label{cycuts:eq:resolution-product}
 u_Sv_T=(-1)^{s|u|}\bigl(m_p(u\boxtimes v)\bigr)_G.
\end{equation}
The total degree of $u_S$ is $|u|+r$.

\begin{proposition}\label{cycuts:prop:resolution-product}
These operations make $K\oplus\mathcal K(C)$ an augmented $E_1$
cut algebra, and the support comparison identifies it with $\Ocut C$
as an algebra.
\end{proposition}
\begin{proof}
For a third term with block length $h$, the two multiplication signs
have exponents $s|u|+h(|u|+|v|)$ and $h|v|+(s+h)|u|$, which agree.
Both coefficient maps are the inverse factorization map on the same
final open. Its coherences give the higher associativity data.

An added cut in the first word retains its insertion position.
An added cut in the second word changes its insertion position by
$r$. The Leibniz sign and changed product sign contribute
$|u|+r+(s+1)|u|=r+s|u|$ modulo two, as required. An internal
differential on the first coefficient gives the equal exponents
$s|u|+r+s$ and $r+s(|u|+1)$; the second coefficient is identical
after adding $|u|$ to both. Thus the differential is a derivation.
Finally, the consolidation exponent for concatenated block lists
differs from the two separate exponents by $s|u|$. Equation
\eqref{cycuts:eq:consolidation} therefore intertwines this product
with cobar concatenation, including the unit and coherent faces.
\end{proof}

Call an augmented cut algebra \emph{$Z$-supported} when its
augmentation ideal in arity $n$ is supported on $Z_n$.

\begin{theorem}\label{cycuts:thm:strong}
Bar and cobar restrict to inverse equivalences between $Z$-supported
augmented cut algebras and strong coaugmented cut coalgebras, with
empty coefficient $k$ on both sides.
\end{theorem}
\begin{proof}
Theorem~\ref{cycuts:thm:support} gives one direction. Conversely,
restrict the bar of $A=K\oplus a$ to a cut open $U_p$. A factor
whose interval crosses $p$ has an open internal cut there, so its
$Z$-supported coefficient vanishes. Each surviving bar word has
a separator at $p$. Deconcatenation identifies it with a word on
the prefix and a word on the suffix, with the same coefficient.
Every multiplication face crossing that separator has zero target;
the other faces and all higher coherences belong to the two factors.
The supplied coproduct component is consequently an equivalence.
The ambient inverse functors of Theorem~\ref{cycuts:thm:duality}
restrict to these full subcategories and their mapping spaces.
\end{proof}

An equality partition is \emph{cut-indecomposable} if each gap is
crossed by one of its blocks. Its exact equality stratum $S_\pi$
then lies in $Z_n$; write $i_\pi:S_\pi\hookrightarrow X^n$ for its
inclusion. The smaller support consisting only of the
small diagonal requires an additional condition.

\begin{corollary}\label{cycuts:cor:small-diagonal}
For a strong cut coalgebra $C$, its reduced cobar is supported on
the small diagonal in every arity if and only if
\[
 i_\pi^!C_n=0
\]
for each cut-indecomposable partition $\pi$ with more than one
equality block. In that case the proper diagonal counit gives
\[
 \Delta_{n,*}(C_1[-1])\simeq\overline{\Ocut C}_n.
\]
\end{corollary}
\begin{proof}
On a stratum in $Z_n$, extraordinary restriction of
$R\Gamma_{Z_n}C_n$ is that of $C_n$; outside $Z_n$ it vanishes.
The finite filtration by equality strata gives the stated criterion.
Kashiwara equivalence and the monotone collapse to the singleton
identify the remaining coefficient and its counit.
\end{proof}

\subsection{A coefficient beyond binary data}

Let $X$ have two distinct geometric points and let $J$ be
\eqref{cycuts:eq:J-word}. Define
\begin{equation}\label{cycuts:eq:hidden}
 P_n=\operatorname{cofib}
 \bigl(R\Gamma_{\Delta_n}J_n\longrightarrow
                       R\Gamma_{Z_n}J_n\bigr)[-1],\qquad P_0=0.
\end{equation}
Both supports pull back correctly under monotone collisions, so $P$
is a word system. It is $Z$-supported. If
$v:X^2\setminus\Delta\hookrightarrow X^2$, then
\[
 P_1=P_2=0,\qquad
 P_3\simeq i_{13,*}v_*J_{X^2\setminus\Delta}[-1].
\]
Set $A^0=\Ocut J$ and give $P$ its $A^0$-bimodule structure through
the augmentation. The square-zero extension $A'=A^0\ltimes P$ has
zero product on $P\cut P$.

\begin{proposition}\label{cycuts:prop:binary-invisible}
The strong cut coalgebras $J$ and $\Bcut A'$ have identical empty,
unary and binary data and are not equivalent.
\end{proposition}
\begin{proof}
Theorem~\ref{cycuts:thm:strong} makes $\Bcut A'$ strong. A positive
bar word containing $P$ needs at least three labels, so projection
to $\Bcut A^0\simeq J$ is an equivalence through arity two.
At arity three the only additional word is $P_3[1]$. Hence
\[
 (\Bcut A')_3\simeq J_3\oplus P_3[1].
\]
On $x_1=x_3\ne x_2$ its extraordinary restriction has two copies
of the dualizing coefficient, while $J_3$ has one. Their ranks differ.
\end{proof}

This example uses derived coefficients. On the affine line, put
$M=i_{13,*}v_*\omega_{X^2\setminus\Delta}$. Then
\[
 (\Bcut A')_3[-3]\simeq\omega_{X^3}\oplus M[-1],
 \qquad H^1((\Bcut A')_3[-3])=M\ne0.
\]
It therefore fails the condition that the arity-$n$ coefficient shifted
by $[-n]$ be an ordinary differential module. A reconstruction theorem
on that smaller class requires its own hypotheses and comparisons.

\section{A collision that destroys cuts}
\label{cycuts:sec:transport}

The preceding support formula is compatible with monotone collisions.
To examine another collision, order the blocks of a set partition by
their last positions. A surjection $\alpha:[n]\twoheadrightarrow[m]$
with this \emph{root order} satisfies
\[
 \max\alpha^{-1}(1)<\cdots<\max\alpha^{-1}(m).
\]
Its retained quotient gaps and source positions are
\begin{equation}\label{cycuts:eq:retained}
 T_\alpha=\{q\in[m-1]:\alpha^{-1}([q])=[p]\text{ for some }p\},
 \qquad p_\alpha(q)=p.
\end{equation}
For these surjections, too, let $d_\alpha:X^m\to X^n$ be the
ordinary diagonal $d_\alpha(y)_i=y_{\alpha(i)}$. Composites retain
root order: the last position of a composite fibre belongs to its
last intermediate fibre. The map
$p_\alpha$ is strictly increasing on its domain.

\begin{proposition}\label{cycuts:prop:retained}
A source cut crossed by an $\alpha$-fibre has empty inverse image.
Every other cut is $p_\alpha(q)$ for a unique $q\in T_\alpha$,
and its inverse image is $U_q^{(m)}$. For root-ordered collisions
$\alpha:[n]\twoheadrightarrow[m]$ and
$\beta:[m]\twoheadrightarrow[\ell]$,
\begin{equation}\label{cycuts:eq:gap-composition}
 T_{\beta\alpha}=
 \{q\in T_\beta:p_\beta(q)\in T_\alpha\},\qquad
 p_{\beta\alpha}(q)=p_\alpha(p_\beta(q)).
\end{equation}
All quotient gaps survive precisely when $\alpha$ is monotone.
\end{proposition}
\begin{proof}
A crossing fibre imposes an equality forbidden by the cut open.
Otherwise each fibre lies on one side. The left fibres have all their
last positions before those of the right fibres, so root order makes
their quotient labels a prefix. This gives the asserted inverse image.

If $\alpha^{-1}(E)$ is a prefix, the maxima of the fibres indexed
by $E$ precede all other maxima; hence $E$ is itself a prefix.
Apply this observation to $E=\beta^{-1}([q])$. It proves
\eqref{cycuts:eq:gap-composition} in both directions. If every
quotient prefix lifts to a prefix, successive differences make all
fibres consecutive intervals in their given order. This is monotonicity.
\end{proof}

For $T\subseteq[m-1]$, put
\[
 Z_T^{(m)}=\bigcap_{q\in T}(X^m\setminus U_q^{(m)}),
 \qquad Z_\varnothing^{(m)}=X^m.
\]
Taking complements in the proposition gives
\[
 d_\alpha^{-1}Z_n=Z_{T_\alpha}^{(m)},\qquad
 Z_m\subseteq Z_{T_\alpha}^{(m)}.
\]
Let $A=\Ocut J$, and write $\bar A$ for its augmentation ideal.
The support inclusion defines a map
\begin{equation}\label{cycuts:eq:Phi}
 \Phi_\alpha:\bar A_m\longrightarrow d_\alpha^!\bar A_n,
 \qquad d_\alpha^!\bar A_n
       \simeq R\Gamma_{Z_{T_\alpha}^{(m)}}J_m[-1].
\end{equation}
In the finite localization cubes it keeps exactly the subsets of
retained gaps. Equation~\eqref{cycuts:eq:gap-composition} and coherent
open base change prove
\[
 \Phi_{\beta\alpha}\simeq d_\beta^!\Phi_\alpha\circ\Phi_\beta.
\]
The map preserves the length-one projection to $J_m[-1]$.
It respects multiplication at separators retained by the collision:
both composites restrict the same inverse factorization map to the
common open, with the signs of \eqref{cycuts:eq:resolution-product}.
No monoidal structure on arbitrary nonmonotone pullback is assumed.

\begin{theorem}\label{cycuts:thm:collision-cone}
Let $O_m=\bigcup_{q=1}^{m-1}U_q^{(m)}$ and $j_m:O_m\hookrightarrow X^m$.
The cone of $\Phi_\alpha$ is
\begin{equation}\label{cycuts:eq:cone}
 E_\alpha\simeq
 j_{m,*}R\Gamma_{Z_{T_\alpha}^{(m)}\cap O_m}(j_m^!J_m)[-1].
\end{equation}
If $X$ has two distinct geometric points, $\Phi_\alpha$ is an
equivalence exactly for monotone $\alpha$.
\end{theorem}
\begin{proof}
Apply the localization triangle for $Z_m$ to
$R\Gamma_{Z_{T_\alpha}^{(m)}}J_m$. Its supported term is
$R\Gamma_{Z_m}J_m$, giving \eqref{cycuts:eq:cone}.

For a missing quotient gap $q$, choose distinct $x,y$ over an
algebraically closed extension $\kappa/k$, and take
$t_q=(x,\ldots,x,y,\ldots,y)$, with $q$ occurrences of $x$.
Only the cut $q$ is open there. The retained-cut cube has just its
empty vertex after extraordinary restriction, giving $\kappa[-1]$.
The full cube is acyclic, because its $q$ direction is an identity.
Consequently
\[
 t_q^!E_{\alpha,\kappa}\simeq\kappa[-1]\ne0.
\]
For a monotone collision every gap is retained, so the two cubes
and their comparison agree.
\end{proof}

For the outer merger $\alpha=(2,1,2)$, no quotient gap survives.
Equation~\eqref{cycuts:eq:Phi} becomes
\begin{equation}\label{cycuts:eq:outer-merger}
 R\Gamma_\Delta J_2[-1]\longrightarrow J_2[-1],
 \qquad E_\alpha=j_*j^!J_2[-1],
\end{equation}
where $j$ includes the distinct pairs. The two endpoint coefficients
already differ on that open. Changing representatives of the map
cannot make them equivalent.

Even a surviving cut retains its full localization map. Take $X=\mathbb A^1$.
For
$\alpha=(2,1,2,3)$, the last source cut has defining equation
$(z_1-z_4)(z_2-z_4)(z_3-z_4)$. With
$x=y_1-y_3$, $y=y_2-y_3$, its pullback is $xy^2$, whereas the
quotient cut is defined by $xy$. For a polynomial numerator $g$
and an integer $a\ge0$, the inverse presentations are
\begin{equation}\label{cycuts:eq:denominators}
 \frac{g}{(xy^2)^a}\longmapsto\frac{x^ag}{(xy)^{2a}},
 \qquad
 \frac{g}{(xy)^a}\longmapsto\frac{y^ag}{(xy^2)^a}.
\end{equation}
They preserve the rational section and commute with differentiation.
The next collision $x=y=d$ sends $1/(xy^2)$ to $d^{-3}$ by either
route. The numerator in the first map is essential: omitting it
would instead give $d^{-4}$.

\section{The ordered collision atlas}
\label{cycuts:sec:atlas}

A word system records monotone identifications. To retain an order
within every collision, start with one copy $X^I_\sigma$ of $X^I$
for each total order $\sigma$ on a finite nonempty set $I$.
For two orders put
\begin{equation}\label{cycuts:eq:overlap}
 V_{\sigma\rho}=X^I\setminus
 \bigcup_{\substack{\{a,b\}\subset I\\
               \sigma,\rho\text{ disagree on }a,b}}\Delta_{ab}.
\end{equation}
Glue the two charts by the identity on this open. On the affine line,
this is the ravioli configuration construction of
\cite{cycuts:AKY}*{\S4.1}. The definition also makes sense for every
separated finite-type scheme, since it uses closed equality loci.
Here $X$ remains smooth of dimension one.

\begin{proposition}\label{cycuts:prop:atlas}
The charts glue to a smooth quasi-compact quasi-separated scheme
$\mathscr R_I$ with a locally isomorphic map $q_I:\mathscr R_I\to X^I$.
Over a geometric tuple with equality blocks $B$, its fibre has
$\prod_B|B|!$ points.
\end{proposition}
\begin{proof}
If two orders disagree on a pair, a third order disagrees with at
least one of them. Hence
$V_{\sigma\rho}\cap V_{\sigma\eta}
 =V_{\sigma\rho}\cap V_{\rho\eta}$, the gluing condition.
There are finitely many noetherian charts and quasi-compact overlaps.
Smoothness is local on them. Two charts represent the same point
over the tuple exactly when their orders agree within each equality
block. Independent total orders on these blocks extend to an order
on $I$, proving the fibre count.
\end{proof}

At two labels the scheme has two diagonal branches and one common
distinct-pair open. At a total triple collision it has six branches.
The chosen branch records the order inside each equality block.

A surjection $\alpha:I\twoheadrightarrow J$ with an order on each
fibre expands a total order $\sigma$ of $J$ into an order
$e_\alpha(\sigma)$ of $I$. Diagonal substitution on those charts
defines
\begin{equation}\label{cycuts:eq:atlas-collision}
 \delta_\alpha:\mathscr R_J\longrightarrow\mathscr R_I,
 \qquad x_i=y_{\alpha(i)}.
\end{equation}
Indeed,
$d_\alpha^{-1}V_{e_\alpha(\sigma),e_\alpha(\rho)}
 =V_{\sigma\rho}$: there is no disagreement inside a fibre, and
disagreements between fibres are exactly the quotient disagreements.
Thus the chart maps glue. Successive expansion concatenates the same
ordered lists, giving
$\delta_\alpha\delta_\beta=\delta_{\beta\alpha}$.

These maps need not be closed immersions. For the ordered fibre
$1<3$ and singleton fibre $2$, the quotient order $2|13$ expands
to chart $213$. Where the quotient roots differ, this chart agrees
with chart $123$ along $x_1=x_3\ne x_2$. At the triple collision,
the image instead lies on branches where $1,3$ are consecutive.
Thus its intersection with a fixed chart need not be the entire
ordinary diagonal of that chart.

Let $\mathsf{OSurj}_+$ denote finite nonempty sets and surjections
with ordered fibres. Define
\begin{equation}\label{cycuts:eq:atlas-category}
 \C=D(k)\times
 \varprojlim_{\mathsf{OSurj}_+}\DMod(\mathscr R_I),
\end{equation}
with transitions $\delta_\alpha^!$. On a nonseparated atlas,
differential modules are obtained by descent from its chart cover.
An object of $\C$ includes all the collision equivalences and
their coherences.

For $I=A\sqcup B$, the locus $U_{A,B}\subset\mathscr R_I$ excluding every cross
equality is canonically the disjoint locus in
$\mathscr R_A\times\mathscr R_B$. Orders on each block remain,
and all their shuffles agree because positions in different blocks
are distinct. With $j_{A,B}$ its inclusion, put
\begin{equation}\label{cycuts:eq:atlas-product}
 (F\star G)_I=
 \bigoplus_{I=A\sqcup B}j_{A,B,*}
              (F_A\boxtimes G_B)|_{U_{A,B}}.
\end{equation}
The sum includes all ordered decompositions, including empty blocks;
an empty block contributes its scalar coefficient.

\begin{theorem}\label{cycuts:thm:atlas-duality}
The product $\star$ is symmetric monoidal on $\C$, with unit $K$
supported at the empty configuration. Augmented $E_1$ algebras and
coaugmented counital $E_1$ coalgebras with empty coefficient $k$ are
exchanged by inverse bar and cobar equivalences for this product.
\end{theorem}
\begin{proof}
A decomposition splitting a collision fibre has empty inverse-image
disjoint locus. Every other decomposition comes from a unique quotient
decomposition. Open base change proves compatibility with the coefficient
systems. All parenthesizations of an iterated product have one summand
for each final ordered partition and the same coefficient on its
cross-disjoint open. This supplies the associativity coherences.
Permuting named blocks gives the symmetry, with the usual Koszul signs;
the empty summands give the unit.

Each positive factor consumes a label. The bar and cobar therefore
have at most $|I|$ reduced factors at $I$. The proof of
Theorem~\ref{cycuts:thm:duality} applies with arbitrary ordered
partitions in place of consecutive partitions: each gap still has
the same localized coefficient on both sides of its identity map.
The finite filtrations and coherent comparisons give the asserted
inverse equivalences.
\end{proof}

Equation~\eqref{cycuts:eq:outer-merger} now has a geometric consequence.
No Cartesian atlas system can have chart coefficients
$\overline{\Ocut J}_2$ and $\overline{\Ocut J}_3$ on the increasing
charts in these two arities. On the distinct quotient roots of the
outer merger, the latter pulls back to $J_2[-1]$, whereas the former
vanishes. The actual atlas map agrees with the ordinary outer diagonal
on this open, so a collision equivalence would identify
zero with a nonzero object. An atlas cobar has its own chart
coefficients; Theorem~\ref{cycuts:thm:atlas-duality} retains that
independent construction.

\section{The dualizing coalgebra and its singular product}
\label{cycuts:sec:dualizing}

On the atlas let $J_I=p_{\mathscr R_I}^!k=\omega_{\mathscr R_I}[|I|]$
and $J_0=k$. Its coproduct for an ordered decomposition is localization
to its disjoint locus. Successive localizations impose the same
inequalities, giving coherent coassociativity. Coaugmentation and
counit are the identity at the empty component and zero elsewhere.
The two decompositions with an empty end give the counit identities.
Define
\[
 \mathscr U=\Omega_\star J.
\]
Its bar is equivalent to $J$ by Theorem~\ref{cycuts:thm:atlas-duality}.
We compute the operation retained by its first coefficients.

Fix a reference order on the labels and local coordinates on $X$.
For $B=\{b_1<\cdots<b_m\}$ let
$c_B=dz_{b_1}\wedge\cdots\wedge dz_{b_m}$. It has degree $-m$
in $J_B$; set $\xi_B=tc_B$, of degree $1-m$. If $B=A\sqcup C$,
define $\epsilon(A,C)$ by $c_A\wedge c_C=\epsilon(A,C)c_B$.
For an ordered partition $P=(B_1|\cdots|B_r)$, let
$j_P:U_P\hookrightarrow\mathscr R_I$ include the locus where different
blocks have disjoint positions. The cobar coefficient at $n=|I|$ is
\begin{equation}\label{cycuts:eq:all-partitions}
 \mathscr U_I=
 \bigoplus_{r=1}^n\ \bigoplus_{I=B_1\sqcup\cdots\sqcup B_r}
              j_{P,*}\omega_{U_P}[n-r],
 \qquad P=(B_1|\cdots|B_r),\quad B_i\ne\varnothing.
\end{equation}
Its multiplication concatenates block lists on the corresponding
disjoint open. Splitting block $B_i=A\sqcup C$ has coefficient
\begin{equation}\label{cycuts:eq:partition-sign}
 (-1)^{\sum_{j<i}(1-|B_j|)+|A|+1}\epsilon(A,C).
\end{equation}
This includes both the suspension and the density permutation.
Thus, writing $a_i=\xi_{\{i\}}$,
\begin{align}
 d\xi_{12}&=a_1a_2-a_2a_1,\label{cycuts:eq:binary-differential}\\
 d\xi_{123}&=a_1\xi_{23}-a_2\xi_{13}+a_3\xi_{12}
             -\xi_{12}a_3+\xi_{13}a_2-\xi_{23}a_1.
             \label{cycuts:eq:ternary-differential}
\end{align}
Every three-factor word occurs twice with opposite signs after the
second differential. In general, splits in different blocks cancel
by the odd derivation rule. For three subblocks of one block, the
two refinements have equal total density sign and opposite suspension
sign. This proves $d^2=0$ in every arity.

On the affine-line chart put $R=k[z_1,z_2]$, $u=z_1-z_2$ and
$L=R[u^{-1}]$. Suppress the common density while retaining its signs.
Put $\xi=\xi_{12}$, $e_{12}=a_1a_2$ and $e_{21}=a_2a_1$.
The binary coefficient is
\begin{equation}\label{cycuts:eq:binary-complex}
 R\xi\xrightarrow{d}Le_{12}\oplus Le_{21},
 \qquad d\xi=e_{12}-e_{21},\qquad |\xi|=-1.
\end{equation}
The map is injective. Its degree-zero cohomology has the explicit
isomorphism
\begin{equation}\label{cycuts:eq:binary-cohomology}
 [(f,g)]\longmapsto(f+g,[f])\in L\oplus L/R,
 \qquad (h,[p])\longmapsto[(p,h-p)].
\end{equation}
Changing a lift of $[p]$ changes the pair by $(r,-r)$ for $r\in R$,
which is a boundary. This proves both well-definedness and inversion.
For a constant right density $\nu$,
$(f\nu)\partial_i=-(\partial_i f)\nu$. All displayed maps
commute with this action.

The regular commutator is a boundary, but
\begin{equation}\label{cycuts:eq:singular}
 \beta=\left[\frac{e_{12}-e_{21}}u\right]
       \longleftrightarrow(0,[u^{-1}])
\end{equation}
is nonzero. A primitive would require $u^{-1}\xi$, which is not
an allowed coefficient in degree $-1$. Normal differentiation generates
all higher principal parts. Restoring densities, the second summand
in \eqref{cycuts:eq:binary-cohomology} is the diagonal direct image
$\Delta_{2,*}\omega_X$.

The coefficient interchange $E(f,g)=(g,f)$ acts in these splitting
coordinates by
\[
 E(h,[p])=(h,[h]-[p]).
\]
The half-difference coordinate $[q]=[(f-g)/2]$ instead gives
$(h,[q])\mapsto(h,-[q])$. A geometric label permutation also
changes coordinates and density: exchanging $z_1,z_2$ sends $u$ to
$-u$ and $dz_1\wedge dz_2$ to its negative.

The full bar retains the homotopy and its multiplication together.
Its reduced binary complex is
\begin{equation}\label{cycuts:eq:binary-bar}
 R\oplus L\oplus L\longrightarrow L\oplus L,
 \qquad(r,f,g)\longmapsto(-r+f,r+g),
\end{equation}
in degrees $-2,-1$. It is surjective with kernel $(r,r,-r)$.
The corresponding cycle is
\[
 [\xi]+[a_1|a_2]-[a_2|a_1].
\]
Its three differential terms cancel by
\eqref{cycuts:eq:binary-differential}. Restoring the density gives
$H^{-2}(B_\star\mathscr U)_2=\omega_{X^2}$ and zero cohomology
in other degrees, hence the chart coefficient is equivalent to
$J_2=\omega_{X^2}[2]$.
The suspended class of $\beta$ is killed by the allowed two-factor
bar word $u^{-1}([a_1|a_2]-[a_2|a_1])$. Thus the bar comparison
retains the singular multiplication and the differential resolving it.

\subsection{Derived coefficients on ordinary position space}

The map $q_I:\mathscr R_I\to X^I$ permits a second coefficient
calculation. Choose an auxiliary ordering of its chart set
$\Sigma_I=\operatorname{Ord}(I)$. For a nonempty subset $S$ of
charts, let $V_S$ remove the diagonals whose pair order varies in
$S$, and let $j_S:V_S\hookrightarrow X^I$. Define the following
complex of left differential modules:
\begin{equation}\label{cycuts:eq:cech}
 \check C_I^p=\bigoplus_{\sigma_0<\cdots<\sigma_p}
       j_{\{\sigma_0,\ldots,\sigma_p\},*}
                   \mathcal O_{V_{\{\sigma_0,\ldots,\sigma_p\}}},
 \qquad d=\sum_a(-1)^a\operatorname{res}_a.
\end{equation}
Different subsets remain different summands even when their opens
coincide.

Write $q_{I,+}$ for the derived direct image of right differential
modules along $q_I$.

\begin{proposition}\label{cycuts:prop:derived-image}
There are equivalences of right differential modules
\[
 q_{I,+}(B_\star\mathscr U)_I\simeq q_{I,+}J_I
 \simeq(\omega_{X^I}\otimes_{\mathcal O_{X^I}}\check C_I)[|I|].
\]
\end{proposition}
\begin{proof}
The first equivalence is the direct image of $J\simeq B_\star\Omega_\star J$.
Since $q_I$ is locally an isomorphism, its transfer module contributes
neither a relative de Rham differential nor a dimension shift. Its
remaining sheaf direct image is computed by the finite chart cover.
Locally on the base, every intersection is localization at a product
of diagonal equations and has affine open inclusion. The augmented
\v Cech resolution is exact, contracted locally by inserting a chart
containing the point. Thus its direct image is
\eqref{cycuts:eq:cech}. The canonical line pulls back from $X^I$;
tensoring by it converts the left connection into a right differential
module and gives the stated shift.
\end{proof}

For two labels,
\[
 \check C_2=[R\oplus R\xrightarrow{(f,g)\mapsto g-f}L],
 \qquad H^0(\check C_2)=R,\quad H^1(\check C_2)=L/R.
\]
The chart coefficient $J_2$ and its derived direct image therefore
have different cohomological descriptions.

A commutative dg model retains the products of the derived functions.
On an affine coordinate open $V\subset X^I$, let
$R_S=\Gamma(V\cap V_S,\mathcal O)$. Polynomial forms on the simplex
with vertices $S$ form the differential graded algebra
\[
 \Omega_{\mathrm{poly}}(\Delta^S)=
 k[t_\sigma,dt_\sigma:\sigma\in S]/
           (\textstyle\sum_\sigma t_\sigma-1,\sum_\sigma dt_\sigma),
\]
where $t_\sigma$ has degree zero and $dt_\sigma$ degree one.
Define $\operatorname{TW}_I(V)$ to consist of families
\[
 \eta_S\in R_S\otimes_k\Omega_{\mathrm{poly}}(\Delta^S),
 \qquad S\ne\varnothing,
\]
whose face restrictions agree with the localization maps for every
inclusion of subsets. Multiplication and differential are componentwise.
Integration over oriented faces gives a connection-preserving
quasi-isomorphism of complexes
$\operatorname{TW}_I(V)\to\Gamma(V,\check C_I)$. These local models
glue by restriction. One proof of the equivalence filters by simplex
skeleta: polynomial forms relative to the boundary of a $p$-simplex
have a single cohomology class in degree $p$, identified with $k$
by integration; the attaching maps are the alternating restrictions.
The filtration is finite. This is the all-face model of
\cite{cycuts:AKY}*{Theorem 5}.

Integration need not preserve a chosen cochain product strictly.
In the two-endpoint path model, a degree-zero path $f(x)=x$ and
the one-form $g=u^{-1}dx$ satisfy
$\int_0^1fg=1/(2u)$, whereas $f(0)\int_0^1g=0$.
The commutative product is retained in $\operatorname{TW}_I$;
its comparison with an endpoint cochain product includes higher maps.

\section{The maps between the two bar--cobar constructions}
\label{cycuts:sec:comparison}

Restrict an atlas system to its increasing chart and monotone collisions:
\[
 r:\C\longrightarrow\W.
\]
Inclusion and projection of consecutive decompositions give natural maps
\begin{equation}\label{cycuts:eq:lax-colax}
 rF\cut rG\xrightarrow{\iota}r(F\star G)
                       \xrightarrow{\pi}rF\cut rG,
 \qquad\pi\iota=1.
\end{equation}
Every iterated inclusion retains the same consecutive decomposition;
every iterated projection selects it. The collision comparisons of
the products show that $\iota$ is lax monoidal and $\pi$ is colax
monoidal. Thus $r$ transports algebras by $\iota$ and coalgebras
by $\pi$; write the resulting functors as $r_{\mathrm{alg}}$ and
$r_{\mathrm{co}}$.

These products have different coefficients. For $n$ positive inputs
with rank-one singleton coefficients at $n$ distinct positions, the
all-subset product has $n!$ summands and the cut product has one.
The retraction in \eqref{cycuts:eq:lax-colax} preserves this difference.

\begin{proposition}\label{cycuts:prop:bar-cobar-maps}
For an augmented atlas algebra $A$ and a coaugmented atlas coalgebra
$C$, with empty coefficient $k$, there are natural maps
\begin{align}
 \Bcut r_{\mathrm{alg}}A&\longrightarrow r_{\mathrm{co}}B_\star A,
                  \label{cycuts:eq:bar-map}\\
 r_{\mathrm{alg}}\Omega_\star C&\longrightarrow\Ocut r_{\mathrm{co}}C.
                  \label{cycuts:eq:cobar-map}
\end{align}
The first is a coalgebra map including consecutive words; the second
is an algebra map projecting onto them. On strict normalized models,
the second is a degreewise surjection with a dg two-sided ideal as kernel.
\end{proposition}
\begin{proof}
Adjacent merges of consecutive blocks remain consecutive. Inclusion
therefore commutes with bar faces, internal differentials and higher
faces. Deconcatenation followed by the cut projection retains exactly
the cuts of that consecutive word, proving the coalgebra assertion.

For cobar, a nonconsecutive partition cannot refine to a consecutive
one: merging two adjacent intervals in such a refinement would make
the original partition consecutive. Hence discarded summands have
only discarded refinements. On retained summands the retained
refinements are exactly the cut coproduct, with the same signs and
localizations. Projection is therefore a chain map. Concatenation of
a prefix partition and a suffix partition is consecutive exactly
when both are, proving multiplicativity and the ideal assertion.
Empty terms are unchanged.

A monotone collision kills a partition splitting one of its fibres.
A surviving partition descends to a quotient partition and is
consecutive precisely when that quotient is. Thus both maps respect
collisions. These constructions commute with the normalized finite
diagrams, including their coherent faces, so they also give the
stated enhanced maps. The strict-model description does not assert
strictification on unchanged coefficient objects.
\end{proof}

Projection in the reverse direction for the bar generally fails to
commute with the differential. A discarded reversed word $[a_2|a_1]$
can have differential $[a_2a_1]$, which lies in the retained one-block
term. In a free atlas algebra that product is nonzero. Thus the
bar map in \eqref{cycuts:eq:bar-map} has an essential direction.

For the dualizing coalgebra, the binary cobar map is especially simple:
\begin{equation}\label{cycuts:eq:binary-quotient}
 \begin{array}{ccc}
 R&\xrightarrow{(1,-1)}&L\oplus L\\
 {\scriptstyle1}\big\downarrow&&
        \big\downarrow{\scriptstyle(f,g)\mapsto f}\\
 R&\xrightarrow{1}&L.
 \end{array}
\end{equation}
The rows have degrees $-1,0$; density bases fix the displayed signs.
Its kernel is $L$ in degree zero. On cohomology it gives
\[
 0\longrightarrow L\xrightarrow{h\mapsto[(0,h)]}
 (L\oplus L)/\{(a,-a):a\in R\}
 \xrightarrow{[(f,g)]\mapsto[f]}L/R\longrightarrow0.
\]
Indeed, if $f\in R$, subtracting $(f,-f)$ leaves the unique
representative $(0,f+g)$ of its kernel class. After inverting $u$,
the full cobar has one nonzero degree-zero coefficient and the cut
cobar is acyclic. This comparison is not an equivalence.

\subsection{The ternary kernel}

For the affine line put $R_n=k[z_1,\ldots,z_n]$, $z_{ij}=z_i-z_j$,
and
\begin{equation}\label{cycuts:eq:localization-modules}
 L_n=R_n\left[\left(\prod_{i<j}z_{ij}\right)^{-1}\right],
 \qquad
 M_P=R_n[z_{ij}^{-1}:i,j\text{ lie in different blocks of }P].
\end{equation}
All modules are tensored with the common right density, suppressed
in the coordinate matrices. Equal rings with different partition
labels remain separate summands. Let $F_n$ be the degreewise kernel of
\eqref{cycuts:eq:cobar-map} on this normalized model for $J$, and put
$e_P=\xi_{B_1}\cdots\xi_{B_r}$. Then
\begin{equation}\label{cycuts:eq:discarded}
 F_n^{r-n}=\bigoplus_{\substack{P\text{ has }r\text{ blocks}\\
                              P\text{ not consecutive}}}M_Pe_P.
\end{equation}
The differential has the signs of \eqref{cycuts:eq:partition-sign}.

At three positions, order the degree-minus-one terms by
$13|2,3|12,23|1,2|13$, and the degree-zero terms by
$132,312,321,231,213$. The matrix is
\begin{equation}\label{cycuts:eq:ternary-kernel}
 D=\begin{pmatrix}
 1&0&0&0\\-1&1&0&0\\0&-1&-1&0\\0&0&1&-1\\0&0&0&1
 \end{pmatrix}:
 \bigoplus_P M_P\longrightarrow L_3^5.
\end{equation}
An entry $\pm1$ denotes a signed localization inclusion. Replace
the target coordinates $(a_0,\ldots,a_4)$ by
\[
 \left(\sum_i a_i,\ a_0,\ a_0+a_1,\ a_0+a_1+a_2,
                                  \ a_0+a_1+a_2+a_3\right).
\]
Successive differences invert this integral basis change. It turns
$D$ into a zero row followed by diagonal entries $1,1,-1,-1$.
Consequently
\begin{equation}\label{cycuts:eq:ternary-cohomology}
 H^{-1}(F_3)=0,\qquad
 H^0(F_3)\simeq\omega_{\mathbb A^3}\otimes_{R_3}
       \left(L_3\oplus\bigoplus_P L_3/M_P\right).
\end{equation}
In particular, the class of $z_{13}^{-1}$ in $L_3/M_{13|2}$ is
nonzero: at the generic point of $\Delta_{13}$, the factors
$z_{12},z_{23}$ are units and every element of $M_{13|2}$ is
regular, whereas $z_{13}^{-1}$ has a pole. Fully localizing every
coefficient before taking cohomology would lose this class.

\section{Outer principal parts in every arity}
\label{cycuts:sec:outer}

The discarded complex has negative cohomology beginning at four
positions. Its first class has an explicit representative in every
arity. In the fully localized complex, let $T_n$ be the part of
$d\xi_{[n]}$ keeping the first and last labels in the same block:
\begin{equation}\label{cycuts:eq:Tn}
 T_n=\sum_{\substack{A\sqcup C=[n],\ A,C\ne\varnothing\\
                         1,n\text{ in the same part}}}
       (-1)^{|A|+1}\epsilon(A,C)\xi_A\xi_C,
 \qquad b_n=dT_n.
\end{equation}
All its terms are discarded. Its degree is $2-n$, and $b_n$ has
degree $3-n$. Write $u=z_{1n}$.

\begin{theorem}\label{cycuts:thm:outer}
For $n\ge3$ there is an injection of right differential modules
\begin{equation}\label{cycuts:eq:outer-injection}
 \omega_{\mathbb A^n}\otimes_{R_n}(R_n[u^{-1}]/R_n)
 \lhook\joinrel\longrightarrow H^{3-n}(F_n),
 \qquad [f]\longmapsto[fb_n].
\end{equation}
Thus these negative cohomology groups have no lower bound independent
of arity.
\end{theorem}
\begin{proof}
The permutahedron is the convex hull of the permutations of
$(1,\ldots,n)$ in $\mathbb R^n$. After extending coefficients to $L_n$,
ordered partitions form its cellular chain complex with degrees negated.
The face of a partition with $r$ blocks has dimension $n-r$;
the density orientations give exactly
\eqref{cycuts:eq:partition-sign}. Convexity makes this cellular
complex acyclic below its degree-zero class. The kernel on its
facets is the rank-one $L_n$-module spanned by the boundary of the unique top cell.
That boundary has coefficient $\pm1$ on every two-block partition.
The discarded facets omit all $n-1$ consecutive two-block partitions,
so the differential on their span is injective. The same holds on
their actual localization submodules.

If a three-block partition still keeps $1,n$ together, both its
facet coarsenings occur in $T_n$ and cancel in $dT_n$. Every nonzero
coordinate of $b_n$ therefore separates $1,n$ and inverts $u$.
Thus $fb_n$ is an allowed cochain for every $f\in R_n[u^{-1}]$,
and it is closed. For $f\in R_n$ its primitive $fT_n$ is allowed.

Conversely, suppose $fb_n=dv$ for an allowed discarded two-block
cochain $v$. Injectivity of the fully localized facet differential
forces $v=fT_n$. A facet of $T_n$ does not invert $u$. Unique
factorization gives $M_P\cap R_n[u^{-1}]=R_n$, so its coordinate
$\pm f$ can belong to $M_P$ only when $f\in R_n$. This proves
injectivity. Such a facet exists for $n\ge3$, and the same
injectivity proves $b_n\ne0$. Its constant coefficients make
differentiation of $fb_n$ equal to differentiation of $f$ followed
by the map. Restoring the density proves differential-module linearity.
\end{proof}

For four positions,
\begin{equation}\label{cycuts:eq:T4}
 T_4=-\xi_{124}\xi_3+\xi_{134}\xi_2-\xi_{14}\xi_{23}
     -\xi_2\xi_{134}-\xi_{23}\xi_{14}+\xi_3\xi_{124}.
\end{equation}
The cycle $z_{14}^{-1}dT_4$ is allowed, but its fully localized
primitive $z_{14}^{-1}T_4$ is not. In particular,
$F_4\to H^0(F_4)[0]$ is not a quasi-isomorphism.

This class also fixes the sign of the connecting map from the cut
cobar. Let
\[
 U_n=d_{\mathrm{cut}}\xi_{[n]}
       =\sum_{p=1}^{n-1}(-1)^{p+1}\xi_{[1,p]}\xi_{[p+1,n]}.
\]
Every term inverts $u$, so $fU_n$ is a cut cocycle for a principal
part $f$. Split the full top boundary as $U_n+T_n+V_n$, where
$V_n$ consists of discarded facets separating the outer labels.
The exact-sequence connecting map
$H^{2-n}((\Ocut J)_n)\to H^{3-n}(F_n)$ then satisfies
\begin{equation}\label{cycuts:eq:connecting}
 [fU_n]\longmapsto-[fb_n].
\end{equation}
Indeed, differentiating its full-complex lift gives
$d(fU_n)=-fb_n-d(fV_n)$, and $fV_n$ is an allowed discarded
cochain. The obstruction belongs to this comparison of cobars;
it leaves both ambient bar--cobar equivalences intact.

\subsection{The complete negative group at four positions}

The four-point complex has eleven discarded facets, thirty-three
discarded edges and twenty-three vertices:
\begin{equation}\label{cycuts:eq:F4}
 \bigoplus_{11}M_P\xrightarrow{A}\bigoplus_{33}M_Q
                      \xrightarrow{B}L_4^{23},
\end{equation}
in degrees $-2,-1,0$. After full localization, the ranks of $A,B$
are $11,22$. The first rank follows from the facet argument above.
For the second, the permutahedron graph with vertex $1234$ removed
is connected: increasing inversion number connects every remaining
permutation to $4321$ without visiting $1234$.

Put $R=R_4$, $L=L_4$, and for a pair $h=\{i,j\}$ set
\[
 L_h=R[z_{ab}^{-1}:\{a,b\}\ne h],\qquad W_h=L/L_h.
\]
Let $I_h$ index discarded facets keeping $h$ together, and let
$J_h$ index three-block partitions keeping $h$ together. The submatrix
$K_h=A_{J_h,I_h}$ has the following dimensions and ranks:
\[
 \begin{array}{c|rrrrrr}
 h&12&13&14&23&24&34\\\hline
 |I_h|&4&5&6&4&5&4\\
 |J_h|&5&6&6&5&6&5\\
 \operatorname{rank}K_h&4&5&5&4&5&4.
 \end{array}
\]
These are constant matrices defined by
\eqref{cycuts:eq:partition-sign}. For $h=14$, ordering the columns
as in \eqref{cycuts:eq:T4} gives
\[
 K_{14}=\begin{pmatrix}
 1&0&-1&0&0&0\\0&1&1&0&0&0\\-1&0&0&1&0&0\\
 0&0&0&-1&1&0\\0&-1&0&0&0&1\\0&0&0&0&-1&-1
 \end{pmatrix}.
\]
Its kernel is generated by $(-1,1,-1,-1,-1,1)$, the coefficients
of $T_4$. For the other pairs, lexicographic ordering of block lists
gives the following row equations:
\[
 \begin{array}{c|l}
 12&-v_1,\ -v_2,\ v_2+v_3,\ v_1-v_4,\ -v_3+v_4\\
 13&-v_1,\ v_1-v_2,\ -v_3,\ v_3+v_4,\ v_2+v_5,\ -v_4-v_5\\
 23&v_1,\ -v_2,\ v_2-v_3,\ -v_1+v_4,\ v_3-v_4\\
 24&v_1,\ v_2,\ -v_1-v_4,\ v_3+v_4,\ -v_2+v_5,\ -v_3-v_5\\
 34&v_1,\ v_2,\ -v_2+v_3,\ -v_1-v_4,\ -v_3+v_4.
 \end{array}
\]
Setting each row list to zero successively kills every coordinate.
These eliminations use only $\pm1$, so the same kernel descriptions
hold with coefficients in $W_h$.

\begin{theorem}\label{cycuts:thm:four-negative}
Globally on the standard four-position chart,
\begin{equation}\label{cycuts:eq:four-negative}
 H^{-2}(F_4)=0,\qquad
 H^{-1}(F_4)\simeq
 \omega_{\mathbb A^4}\otimes_R(R[z_{14}^{-1}]/R).
\end{equation}
The second map is $[f]\mapsto[fdT_4]$.
\end{theorem}
\begin{proof}
The facet differential is injective, proving the first assertion.
For a degree-minus-one cocycle $w$, full localization gives a unique
$v\in L^{11}$ with $Av=w$. A three-block term keeping $h$
together has coefficient $L_h$. Thus $w$ belongs to its actual
coefficient modules precisely when
\[
 K_h(v_i\bmod L_h)_{i\in I_h}=0\quad\text{for every }h.
\]
For $h\ne14$, the displayed eliminations force the corresponding
coordinates of $v$ to be regular along $h$.

Unique factorization gives
$M_P=\bigcap_{h\text{ internal to }P}L_h$ inside $L$.
Writing $u=z_{14}$, it follows that $v_P\in M_P[u^{-1}]$.
The coordinates of facets separating $1,4$ already lie in $M_P$
and can be subtracted as legal boundaries. Modulo those boundaries,
the other six principal parts are proportional to the coefficients
of $T_4$. The remaining module is therefore the intersection, in
$W_{14}$, of the six images $M_P[u^{-1}]/M_P$ for $P\in I_{14}$.

Use coordinates $u=z_{14}$, $a=z_{12}$, $b=z_{23}$, $y=z_4$,
and write $c=a+b=z_{13}$, $D=k[y,a,b]$. Finite negative Laurent
parts at $u=0$ identify
\[
 W_{14}\simeq\bigoplus_{j\ge1}D[(abc)^{-1}]u^{-j}.
\]
The factors $u-a,u-c$ need only finitely many Taylor coefficients
for any fixed pole order. Vanishing of the negative part means
divisibility by the required power of $u$, hence membership in $L_{14}$.
For $124|3$, $134|2$ and $14|23$, the images are respectively
\[
 \bigoplus_{j\ge1}D[(bc)^{-1}]u^{-j},\qquad
 \bigoplus_{j\ge1}D[(ab)^{-1}]u^{-j},\qquad
 \bigoplus_{j\ge1}D[(ac)^{-1}]u^{-j}.
\]
Reversed facets repeat these three modules. Since $a,b,c=a+b$
are pairwise nonassociated prime factors in $D$, the intersection
of their coefficient rings is $D$. The desired intersection is
therefore $\bigoplus_{j\ge1}Du^{-j}=R[u^{-1}]/R$.
This identifies the class with $[fdT_4]$, and the density argument
in Theorem~\ref{cycuts:thm:outer} proves differential-module linearity.
\end{proof}

The degree-zero cohomology also has a global finite presentation.
Let $G$ be the permutahedron graph on the orders of $[4]$ other than $1234$,
fix $\sigma_0=4321$, and label each edge by its swapped pair $h$.
Set
\begin{equation}\label{cycuts:eq:graph-module}
 \mathcal G_R=
 \frac{\displaystyle\bigoplus_{\sigma\in V(G)\setminus\{\sigma_0\}}L e_\sigma}
 {\displaystyle\sum_{\{\sigma,\tau\}\in E(G)}
                   L_{h(\sigma,\tau)}(e_\sigma-e_\tau)},
 \qquad e_{\sigma_0}=0.
\end{equation}
Then
\begin{equation}\label{cycuts:eq:four-zero}
 H^0(F_4)\simeq\omega_{\mathbb A^4}\otimes_R(L\oplus\mathcal G_R).
\end{equation}
To prove this, sum all vertex coordinates. Every incidence column
has sum zero, and the $\sigma_0$ coordinate splits this sum.
In its kernel use $e_\sigma-e_{\sigma_0}$ as basis. The edge
relations are precisely \eqref{cycuts:eq:graph-module}, with twenty-two
generator modules and thirty-three relation modules. Constant basis
changes preserve the differential-module action. This presentation
retains the extensions across intersections of the six diagonals;
a decomposition into indecomposable modules is additional information.

\subsection{The crossed collision}

On the open set where the blocks $\{1,3\}$ and $\{2,4\}$ are
separated, write
\[
 \begin{gathered}
 u=z_{13},\quad v=z_{24},\quad
 A=R[(z_{12}z_{14}z_{23}z_{34})^{-1}],\\
 A_u=A[u^{-1}],\quad A_v=A[v^{-1}],\quad L=A[(uv)^{-1}].
 \end{gathered}
\]
Here $u,v$ are independent normal coordinates. Suppressing the common
density, the coefficient types of $F_4$ are
\begin{equation}\label{cycuts:eq:crossed-types}
 A^2\oplus A_u^3\oplus A_v^3\oplus L^3
 \longrightarrow A_u^6\oplus A_v^6\oplus L^{21}
 \longrightarrow L^{23}.
\end{equation}

A scalar pivot between two \emph{identical} coefficient modules
can be cancelled by a chain homotopy. For a nonzero scalar pivot
$a=D_{ij}$ in row $i$, column $j$ of a differential matrix $D$, the
remaining matrix is the Schur complement
\begin{equation}\label{cycuts:eq:schur}
 D'_{rc}=D_{rc}-D_{rj}a^{-1}D_{ic}.
\end{equation}
Triangular row and column operations isolate the two-term isomorphism,
whose contracting homotopy is $a^{-1}$. All other entries remain
composites of allowed localization maps. The equality of the pivot
modules is necessary: a localization inclusion is not an isomorphism
merely because its displayed scalar is $1$.

Order partitions lexicographically by their increasing block lists.
Choose the first nonzero pivot between identical modules, first by
lower degree, then row and column; retain the order after cancellation.
The nine pivots from degree $-2$ have types
\[
 A_u,L,A_v,A_v,L,A_u,A_u,L,A_v
\]
and scalars $1,1,-1,-1,1,1,1,1,1$. The next eighteen pivots have
type $L$ and scalar $1$. Formula~\eqref{cycuts:eq:schur} leaves
\[
 A^2\xrightarrow{M_0}A_u^3\oplus A_v^3\xrightarrow{M_1}L^5,
\]
with retained middle order $(A_u,A_v,A_v,A_u,A_u,A_v)$ and matrices
\begin{equation}\label{cycuts:eq:crossed-matrices}
 M_0=\begin{pmatrix}1&0\\-1&0\\0&1\\0&-1\\-1&1\\1&-1\end{pmatrix},
 \qquad
 M_1=\begin{pmatrix}
 1&1&0&0&0&0\\0&0&1&1&0&0\\0&-1&-1&0&1&0\\
 -1&0&0&-1&0&1\\0&0&0&0&-1&-1
 \end{pmatrix}.
\end{equation}
For completeness, the required final basis changes are
\[
 S=\begin{pmatrix}
 1&0&1&0&0&0\\0&0&0&1&0&1\\0&0&0&0&-1&0\\
 0&-1&0&0&0&0\\-1&1&0&0&0&0\\0&0&0&-1&1&0
 \end{pmatrix},\qquad
 W=\begin{pmatrix}
 1&0&1&1&1\\0&-1&0&0&0\\-1&1&0&-1&0\\
 -1&1&-1&0&0\\1&-1&0&0&0
 \end{pmatrix}.
\]
Their determinants are units. The first three columns of $S$ have
type $A_u$ and the last three type $A_v$, so it respects the
coefficient modules. Direct multiplication gives
\begin{equation}\label{cycuts:eq:crossed-normal-form}
 S^{-1}M_0=
 \begin{pmatrix}1&0\\0&1\\0&0\\-1&0\\0&-1\\0&0\end{pmatrix},
 \qquad
 W^{-1}M_1S=
 \begin{pmatrix}
 1&0&0&1&0&0\\0&1&0&0&1&0\\0&0&1&0&0&0\\
 0&0&0&0&0&1\\0&0&0&0&0&0
 \end{pmatrix}.
\end{equation}

\begin{theorem}\label{cycuts:thm:crossed}
On this open set, $F_4$ is chain-homotopy equivalent by
differential-module maps to
\begin{equation}\label{cycuts:eq:crossed-splitting}
 L[0]\oplus[A_u\to L]\oplus[A_v\to L]
       \oplus[A\to A_u\oplus A_v\to L]^{\oplus2}.
\end{equation}
The two-term complexes have degrees $-1,0$ and the three-term
complexes degrees $-2,-1,0$. Hence
\[
 H^{<0}(F_4)=0,\qquad
 H^0(F_4)\simeq\omega\otimes_A
 \left(L\oplus L/A_u\oplus L/A_v
                \oplus(L/(A_u+A_v))^{\oplus2}\right).
\]
\end{theorem}
\begin{proof}
The two paired columns in \eqref{cycuts:eq:crossed-normal-form}
give complexes with maps $a\mapsto(a,-a)$ and $(f,g)\mapsto f+g$.
The third and sixth columns give the two single localization maps;
the last target row gives $L[0]$. The preceding cancellations and
constant basis changes prove \eqref{cycuts:eq:crossed-splitting}.

The localization maps are injective. In either three-term complex,
the middle kernel consists of $(f,-f)$ with $f\in A_u\cap A_v$.
Unique factorization and independence of $u,v$ give $A_u\cap A_v=A$.
Only the displayed cokernels therefore remain. All transformations
are scalar matrices or localization maps, so they commute with the
right differential operators.
\end{proof}

Both codimension-two summands are nonzero. Near $u=v=0$, the
double residue of $1/(uv)$ is $1$, while it vanishes on sections
having poles in only one of these variables. The cross-block
denominators are units there. Thus $1/(uv)$ survives modulo
$A_u+A_v$. On this open $z_{14}$ is invertible, so the negative
outer module of \eqref{cycuts:eq:four-negative} vanishes, consistently
with the local calculation.

\section{Exchange and the two factorization conditions}
\label{cycuts:sec:exchange}

To recover an atlas system from a word system, one must specify how
its charts agree on their full overlaps. For a word system $F$,
let $F_I^\sigma$ be its transport to the chart with order $\sigma$.
An \emph{exchange datum} consists of equivalences
\[
 g_{\sigma\rho}:F_I^\sigma|_{V_{\sigma\rho}}
                  \xrightarrow{\sim}F_I^\rho|_{V_{\sigma\rho}},
\]
with coherent \v Cech cocycles and relabeling compatibility. It also
requires collision compatibility: expanding an ordered-fibre
surjection into charts must intertwine these maps with the monotone
collision equivalences, coherently under every composite.
The overlaps in \eqref{cycuts:eq:overlap} may contain partial collisions.

\begin{proposition}\label{cycuts:prop:descent}
Restriction identifies $\C$ with word systems equipped with this
exchange datum. The corresponding tensor is the all-subset product
glued chartwise. Componentwise strong atlas coalgebras correspond
to strong cut coalgebras with factorization-compatible exchange.
\end{proposition}
\begin{proof}
Differential modules satisfy descent for the finite open chart cover.
The $g$ maps and their cocycles therefore glue $F_I^\sigma$ to a
coefficient on $\mathscr R_I$, and descent of morphisms gives the
same statement on mapping spaces. Collision compatibility glues
the chartwise $!$-pullback equivalences and their higher coherences.
Restriction and gluing are inverse.

For an arbitrary ordered decomposition, choose a chart in which its
blocks are consecutive. Its coefficient product is then the cut
product on that chart. The exchange compatibilities glue these
products to \eqref{cycuts:eq:atlas-product}. The same choice identifies
each supplied factorization component with a cut component, proving
the last assertion.
\end{proof}

Write $\operatorname{Conf}_2(X)=X^2\setminus\Delta$ for distinct pairs.
Such data need not exist for an underlying word system. Suppose an
augmented cut algebra has augmentation ideal $a_n=\Delta_{n,*}M$,
and $(M\boxtimes M)|_{\operatorname{Conf}_2(X)}\ne0$. On distinct
pairs its cut bar has coefficient $(M\boxtimes M)[2]$. On the
stratum $x_1=x_3\ne x_2$ its ternary coefficient has zero
extraordinary restriction: a one-block word requires the small
diagonal, a two-block word requires $12$ or $23$ to coincide,
and a three-block word requires all positions to be distinct.
An outer-merger collision equivalence in the atlas would identify
this zero coefficient with the nonzero binary one. Thus the cut bar
has no extension with unchanged chart coefficients.

\subsection{Joint factorization}

For the atlas let $D_I$ be the union of the small-diagonal branches
$\delta_\sigma:X\hookrightarrow\mathscr R_I$. A coefficient
$N\in\DMod(X)$ has the diagonal insertion
\[
 (i_{\mathrm{ord}}N)_I=\bigoplus_\sigma\delta_{\sigma,*}N,
 \qquad(i_{\mathrm{ord}}N)_0=0.
\]
An ordered-fibre collision pulls back exactly the branches whose
orders are expanded quotient orders. Proper base change gives the
collision structure; Kashiwara equivalence identifies every positive
system supported on $D_I$ with $i_{\mathrm{ord}}$ of its singleton.

For an atlas coalgebra $C$ put $N=C_1$ and let $U_I$ be its
fully distinct locus. Its iterated coproducts define a joint map
\begin{equation}\label{cycuts:eq:joint}
 \Theta_I:C_I|_{U_I}\longrightarrow
 \bigoplus_{\sigma\in\operatorname{Ord}(I)}
       (N_{\sigma(1)}\boxtimes\cdots\boxtimes N_{\sigma(n)})|_{U_I}.
\end{equation}
Call $C$ \emph{jointly factorizing} if these supplied maps are
equivalences in every arity.

\begin{theorem}\label{cycuts:thm:joint}
Atlas bar and cobar restrict to inverse equivalences between augmented
algebras with positive coefficients supported on $D_I$ and jointly
factorizing coalgebras. For such a coalgebra,
\[
 i_{\mathrm{ord}}(C_1[-1])\simeq\overline{\Omega_\star C}.
\]
\end{theorem}
\begin{proof}
On distinct positions, a diagonal-supported bar factor occupying
more than one label vanishes. Its surviving terms are the singleton
words, one per total order, and every multiplication face has zero
target. Their deconcatenation is exactly \eqref{cycuts:eq:joint}.

Conversely, projection to singletons followed by all iterated
coproducts maps the separated coefficient system of $C$ to the
tensor coalgebra on its singleton. Its component is $\Theta_I$,
so joint factorization makes it an equivalence. The gap contraction
of Theorem~\ref{cycuts:thm:duality} contracts its cobar at more
than one distinct position, leaving $C_1[-1]$ at one position.
Every equality stratum with at least two roots is the distinct
locus of an ordered-fibre collision. Collision compatibility
therefore kills its extraordinary restriction. Finite-stratum
recollement puts the reduced cobar on $D_I$, and the proper branch
counits give the displayed equivalence. The ambient bar--cobar
maps restrict to these full subcategories and their mapping spaces.
\end{proof}

Now assume that \emph{each} binary coproduct component is an
equivalence on distinct pairs. Using one of them identifies the
joint map with
\[
 (1,R):V\longrightarrow V\oplus V,
 \qquad V=(N\boxtimes N)|_{\operatorname{Conf}_2(X)},
\]
where $R$ is an equivalence. Its cofiber is $V$, even for an
arbitrary chain endomorphism $R$. Explicitly, cone coordinates
$(x,y,z)\in V^q\oplus V^q\oplus V^{q+1}$ have differential
\[
 D(x,y,z)=(dx+z,dy+Rz,-dz).
\]
The maps $q(x,y,z)=y-Rx$, $s(v)=(0,v,0)$ and
$H(x,y,z)=(0,0,x)$ satisfy $qs=1$ and $DH+HD=1-sq$.
Thus the ordinary graph kernel is zero but its derived fibre is
$V[-1]$, and the separated binary cobar is $V[-2]$.

\begin{corollary}\label{cycuts:cor:intersection}
A componentwise strong atlas coalgebra $C$ is jointly factorizing
if and only if its singleton satisfies
\[
 (C_1\boxtimes C_1)|_{\operatorname{Conf}_2(X)}=0.
\]
In this case its positive coefficients are diagonal-supported and
its reduced coproduct vanishes.
\end{corollary}
\begin{proof}
The graph cofiber proves necessity. Conversely, the assumed vanishing
kills every separated singleton tensor of length at least two.
Successive strong components identify the coalgebra with these tensors
on distinct loci. Ordered-fibre collisions and finite-stratum
recollement put its positive coefficients on the small-diagonal
branches. Their two-factor convolution is zero by the same separated
square condition, so the reduced coproduct vanishes. For $|I|>1$
both sides of \eqref{cycuts:eq:joint} vanish; for $|I|=1$ the map
is the identity.
\end{proof}

For an atlas algebra with diagonal-supported augmentation ideal,
one can also see the component coproduct directly in its bar.
On a disjoint decomposition $I=P\sqcup Q$, every surviving block
has one of these two colors. A multiplication face between colors
has zero target, so words split into subcomplexes indexed by their
maximal consecutive runs of colors. The coproduct for $P|Q$
projects onto the sector with all $P$ blocks before all $Q$ blocks.
That projection has a section as a coefficient complex. Reverse
and alternating sectors remain in its kernel.

Finally, derived primitives have a precise associative description.
For a positive object $V$, let $\operatorname{Prim}_0(V)=K\oplus V$
have zero reduced coproduct and trivial higher reduced cooperations.
In either carrier, bar--cobar duality and the free-algebra property give
\[
 \begin{aligned}
 \operatorname{Map}_{\mathrm{Coalg}_{E_1}}
       (\operatorname{Prim}_0(V),C)
 &\simeq\operatorname{Map}_{\mathrm{Alg}_{E_1}}
       (T(V[-1]),\Omega C)\\
 &\simeq\operatorname{Map}(V,\overline{\Omega C}[1]).
 \end{aligned}
\]
The tensor algebra here uses the chosen geometric product, with
finite words at each configuration. The right adjoint thus retains
derived fibres; an ordinary kernel of the reduced coproduct would
lose them.

\section{Operations on the small diagonals}
\label{cycuts:sec:operations}

The support of a coefficient determines which associativity homotopies
it can carry. In this section let $X$ be a smooth separated scheme
of pure dimension $d\ge1$. The word category and cut product have
the same definitions. For $M\in\DMod(X)$, let
\[
 (iM)_n=\Delta_{n,*}M\quad(n>0),\qquad(iM)_0=0,
\]
where $\Delta_n$ is the small diagonal. Kashiwara equivalence and
monotone collapse identify these coefficients coherently.
Let $j_n:U_n\hookrightarrow X^n$ include the pairwise-distinct locus.
The derived complex of $n$-input chiral operations is
\begin{equation}\label{cycuts:eq:operations}
 P_n(M_1,\ldots,M_n;P)=
 \RHom_{\mathcal D_{X^n}}
 \bigl(j_{n,*}j_n^*(M_1\boxtimes\cdots\boxtimes M_n),
                                  \Delta_{n,*}P\bigr).
\end{equation}
It retains all degrees of the operation, not just ordinary module maps.

\begin{proposition}\label{cycuts:prop:operations}
For every word system $G$, restriction to the singleton-block summand
at $n$ positions gives
\[
 \RHom_{\W}(iM_1\cut\cdots\cut iM_n,G)
 \simeq\RHom_{\mathcal D_{X^n}}
 \bigl(j_{n,*}j_n^*(\boxtimes_iM_i),G_n\bigr).
\]
For $G=iP$, this identifies the mapping complex with
\eqref{cycuts:eq:operations}, compatibly with composition along
monotone partitions.
\end{proposition}
\begin{proof}
Put $Q=j_{n,*}j_n^*(\boxtimes_iM_i)$. At $m$ positions the cut
product is the direct sum of $d_{\alpha,*}Q$ over monotone
surjections $\alpha:[m]\twoheadrightarrow[n]$. Adjunction and
the structure maps of $G$ identify each component mapping complex
with $Y=\RHom(Q,G_n)$. Under another monotone collision, a summand
survives precisely when its partition descends to the quotient;
otherwise its extraordinary restriction meets a forbidden diagonal
of $Q$ and vanishes. The mapping diagram is thus the product of
copies of $Y$ indexed by surjections to $[n]$.

The category of these surjections has final object
$\operatorname{id}_{[n]}$, so its nerve is contractible. The limit
of the mapping diagram is $Y$, detected on that singleton summand.
For composition, an inner operation localizes within each consecutive
block and an outer operation localizes between the blocks. Open
base change identifies the common source with the fully localized
input. The resulting composite is the external product of the
inner maps followed by the direct image of the outer map, with
the Koszul signs for tensoring maps. This is exactly composition
of the displayed chiral operations.
\end{proof}

Now let every $M_i$ and $P$ be an ordinary right differential module.
The object $j_{n,*}j_n^*(\boxtimes_iM_i)$ has cohomology in
degrees $0$ through $(n-1)(d-1)$. To prove the upper bound,
add the points one at a time. Over the distinct first $n-1$
points, the forbidden graphs for the last point are disjoint
smooth loci of codimension $d$. Locally their complement is
covered by $d$ affine opens, giving direct-image amplitude at
most $d-1$. Induction gives the stated bound, also after any
smooth auxiliary factor.

Set $r=(n-1)(d-1)$. Shifting all inputs and the output by $[d-1]$
changes the operation complex to
\[
 \RHom(j_{n,*}j_n^*(\boxtimes_iM_i),\Delta_{n,*}P)[-r].
\]
It has no negative cohomology, and its degree-zero part is
$\operatorname{Hom}(H^r(j_{n,*}j_n^*(\boxtimes_iM_i)),\Delta_{n,*}P)$.
The mapping spaces of operations on $i(M[d-1])$ are therefore
discrete. An associative structure on it is determined by a binary
degree-zero operation whose two ternary composites are equal.

\begin{theorem}\label{cycuts:thm:diagonal-nilpotence}
Every such associative structure on $i(M[d-1])$ has zero ternary
composites. More generally, a degree-zero ternary operation on shifted
modules that factors through both channels $12|3$ and $1|23$ is zero,
even when the intermediate modules differ.
\end{theorem}
\begin{proof}
First let $d=1$ and use an \'etale coordinate near the small diagonal.
Put $s=z_1-z_2$, $v=z_2-z_3$, $h=s+v=z_1-z_3$. The intermediate
module for the first channel has locally nilpotent $s$ and invertible
$h$; for the second it has locally nilpotent $v$ and invertible $h$.
Their first maps are linear over $\mathcal O[h^{-1}]$, and their
last maps to the small-diagonal target are $\mathcal O$-linear.

For a fixed input $e$, choose exponents $a,b$ killing its images
by $s^a,v^b$ in the two intermediate modules. For $N=a+b-1$,
\[
 e=h^{-N}(s+v)^Ne
   =\sum_{j=0}^N\binom Nj s^jv^{N-j}h^{-N}e.
\]
If $j\ge a$, its common output vanishes by the first factorization;
otherwise $N-j\ge b$ and it vanishes by the second. The inverse
of $h$ is used before the final map. This proves the curve case
without a uniform nilpotence or pole bound.

For $d\ge2$, put $F=M_1\boxtimes M_2\boxtimes M_3$ and
$E=H^{2d-2}(j_{3,*}j_3^*F)$. The degree bound identifies a ternary
operation with a module map $E\to\Delta_{3,*}P$. If $Z$ is
the union of the three pair diagonals, recollement gives
$E=H_Z^{2d-1}(F)$. Each pair diagonal has local-cohomology
amplitude at most $d$, and their common small diagonal has
amplitude at most $2d$.

For pair diagonals $A,B$, let $S_{AB}\subset E$ be the image
of $H_{A\cup B}^{2d-1}(F)$. Mayer--Vietoris gives
$H_{A\cup B}^{2d-1}(F)\simeq H_\Delta^{2d}(F)$ and
$H_{A\cup B}^{2d}(F)=0$. With $C$ the third pair diagonal,
the next Mayer--Vietoris sequence is
\[
 0\longrightarrow S_{AB}\longrightarrow E
             \longrightarrow H_\Delta^{2d}(F)\longrightarrow0.
\]
The first map is injective because the preceding map from
$H_\Delta^{2d-1}(F)$ factors through $H_A^{2d-1}(F)=0$.
Naturality for $A\subset A\cup B$ identifies the image of
$S_{AC}$ in the quotient with the connecting isomorphism for
$A,C$. Hence $S_{AB}+S_{AC}=E$.

An inner operation through $12|3$ is a map
$j_*j^*(M_1\boxtimes M_2)\to\Delta_*N[1-d]$.
Its restriction to $M_1\boxtimes M_2$ is a negative Ext class
between modules and is zero. The composition formula of
Proposition~\ref{cycuts:prop:operations} therefore makes the ternary
operation vanish on $S_{AB}$ for $A=\Delta_{13}$ and
$B=\Delta_{23}$. A factorization through $1|23$ makes it vanish
on $S_{AC}$ for $C=\Delta_{12}$. Their sum is $E$, proving
the assertion. Associativity identifies these two factorizations.
Every binary tree with at least three leaves contains two adjacent
internal vertices, so all longer composites vanish as well.
\end{proof}

The hypothesis is a shifted ordinary module on the small diagonals.
The cut-support complexes and atlas cobars above have larger supports
and additional cochain degrees. Their multiplication homotopies and
outer classes occupy precisely the information that this hypothesis
excludes. The finite positive-word bound places no two-step-nilpotence
condition on those larger algebras.

\begin{problem}\label{cycuts:prob:field-comparison}
Construct collision-compatible maps from the oriented Hall
correspondence coefficients to the operation complexes of an
independently specified field theory. The maps must preserve the
coefficient shifts, translations, supported operations and all
substitutions, including the nonconsecutive collisions. If a word
model is used, supply its full-overlap exchange datum and compare
it with the atlas through \eqref{cycuts:eq:bar-map} and
\eqref{cycuts:eq:cobar-map}, retaining the action of their kernel
classes. Specify separately the empty-support unit, the local
vacuum and any prepared state.
\end{problem}

An order in the collision atlas records branches over coincident
positions. A comparison with real-time ordered observables additionally
requires the spacetime operation maps and their support conditions.
The two geometric dualities, their support theorems and their explicit
comparison maps provide coefficient constructions for that problem.

\endgroup
